\documentclass[aspectratio=169,11pt]{beamer}

%% Shared preamble for Programming and Quantitative Methods in Finance Beamer presentations
%% Adapted from SPM-PEM/spmcourse-beamer.sty (English localisation)

\usepackage{xcolor}
\usepackage{amsmath,amssymb,bm,mathtools}
\usepackage{booktabs}
\usepackage{hyperref}
\usepackage[english]{babel}

\setbeamertemplate{footline}[frame number]
\setbeamertemplate{itemize items}[circle]
\setbeamercolor{structure}{fg=red!30!green!60!black}
\setbeamercolor{item}{fg=red!30!green!60!black}
\setbeamercolor{itemize item}{fg=red!30!green!60!black}

\definecolor{slidebg0}{RGB}{255,255,255}
\definecolor{slidebg1}{RGB}{220,220,220}
\definecolor{slidebg2}{RGB}{220,235,255}
\definecolor{slidebg3}{RGB}{220,255,230}
\definecolor{slidebg4}{RGB}{255,235,210}
\definecolor{slidebg5}{RGB}{235,220,255}
\definecolor{slidebg6}{RGB}{255,220,205}

\newcounter{slidebg}
\setcounter{slidebg}{1}
\newcommand{\alternateslidebg}{
  \ifnum\value{slidebg}=1
    \setbeamercolor{background canvas}{bg=slidebg1}\setcounter{slidebg}{2}
  \else\ifnum\value{slidebg}=2
    \setbeamercolor{background canvas}{bg=slidebg2}\setcounter{slidebg}{3}
  \else\ifnum\value{slidebg}=3
    \setbeamercolor{background canvas}{bg=slidebg3}\setcounter{slidebg}{4}
  \else\ifnum\value{slidebg}=4
    \setbeamercolor{background canvas}{bg=slidebg4}\setcounter{slidebg}{5}
  \else\ifnum\value{slidebg}=5
    \setbeamercolor{background canvas}{bg=slidebg5}\setcounter{slidebg}{6}
  \else
    \setbeamercolor{background canvas}{bg=slidebg6}\setcounter{slidebg}{1}
  \fi\fi\fi\fi\fi
}

\ExplSyntaxOn
\NewExpandableDocumentCommand{\tocsectionbg}{m}
  { slidebg\int_eval:n { \int_mod:nn { #1 } { 6 } + 1 } }
\ExplSyntaxOff

\setbeamertemplate{section in toc}{%
  \setbeamercolor{section in toc}{fg=black,bg=\tocsectionbg{\inserttocsectionnumber}}%
  \begin{beamercolorbox}[wd=\textwidth,sep=0.7ex,leftskip=1em,rightskip=1em]{section in toc}
    \inserttocsectionnumber.\enspace\inserttocsection
  \end{beamercolorbox}%
  \vspace{0.3ex}%
}
\newcommand{\newsection}[1]{\begin{frame}[plain]
\centering\vfill
{\usebeamercolor[fg]{structure}\LARGE\bfseries \insertsectionnumber\ #1}
\vfill\end{frame}}
\NewDocumentCommand{\bgsection}{s m}{%
  \alternateslidebg
  \section{#2}%
  \IfBooleanF{#1}{\newsection{#2}}%
}
\newcommand{\titleandoutline}{
  \alternateslidebg
  \frame{\titlepage}
  \begin{frame}{Outline}
  \tableofcontents[hideallsubsections]
  \end{frame}
}

\title{Programming and Quantitative Methods in Finance}
\subtitle{Králik Balázs (Office Hours Wed 15:30-17:00 E275) \\
Justin Huang (Office Hours Thu 11:30-13:00 E275)}
\author{}
\date{}

\begin{document}

\titleandoutline

\bgsection{Course structure}

\begin{frame}{Course structure}
  \begin{itemize}
    \item 24 classes in total, each 90 minutes.
    \item First half: 12 classes on Python programming (Králik Balázs).
    \item Second half: 12 classes on statistics (Justin Huang).
  \end{itemize}
\end{frame}


\begin{frame}{Assessment at a glance}
  \begin{itemize}
    \item Two midterms altogether: one in Python (week 6) and one in statistics (week 12). Each midterm is worth 30\%. Midterms: $2 \times 30\% = 60\%$ of the final grade.
    \item 10 quizzes, five in each half. each quiz is 4\% of the final grade. Quizzes: $10 \times 4\% = 40\%$ of the final grade.
    \item A quiz or midterm counts only when completed in class, closed book.
    \item There is one optional combined make-up exam in the exam period. If you decide to take the make-up exam, it will override all previous scores for the final grade calculation, even if it comes out worse. However, the make-up exam will not cause you to fail, if your previous scores were sufficient to pass.
    \item The python and statistics halves must each be passed separately in order to pass the course.
  \end{itemize}
\end{frame}

\bgsection{Python half}

\begin{frame}{How the Python classes work}
  \begin{itemize}
    \item The first part introduces the week's concepts through explanations
          and examples.
    \item The second part is devoted to independent in-class problem solving.
    \item The emphasis is practical: students learn to find and fix syntax,
          implementation, and logical errors while working.
    \item Questions and productive collaboration are encouraged, provided that
          other students can work undisturbed.
  \end{itemize}
\end{frame}

\begin{frame}{Python half: 12 classes}
  \small
  \begin{tabular}{cll}
    \toprule
    Class & Quiz & Main topic \\
    \midrule
    1 &           & Basics, variables, and strings \\
    2 & Strings   & Numeric types, conversion, booleans, input, and conditionals \\
    3 &           & Sequence types, lists, and loops \\
    4 & Sequences Loops & Functions \\
    5 &           & Search and sorting algorithms ?? \\
    6 & ??        & Python midterm 1 \\
    7 &           & Modules, dates, and the \texttt{math} module \\
    8 & Modules Dates Math          & Pandas I \\
    9 &           & Pandas II \\
    10 & Pandas   & Data visualization \\
    11 &          & Practice and preparation for midterm 2 \\
    12 &          & Python midterm  \\
    \bottomrule
  \end{tabular}
\end{frame}

\end{document}
